\chapter{The Lie algebra of a Lie group}
\label{chap:liealgebra}

The Lie algebra $\mathbf{L}(G)$ of a Lie group $G$ is $T_1 G$ with the bracket
$[x, y] = [x_l, y_l](1)$ of the left invariant vector fields with values $x$, $y$ at the
identity \cite[Definition 5.1.6]{GN26}. We compute it in the chart $\varphi$ of $G$ at the
identity, which identifies $T_1 G$ with the model space $E$, as in \cite[(5.5)]{GN26}. This only
needs the multiplication of $G$ near the identity, read in the chart.

\section{Local group laws}

\begin{definition}[Local group laws]
  \label{def:local-group-law}
  \lean{LocallyConvex.LieGroups.LocalGroupLaw}
  \leanok
  A \emph{local group law} on $E$ is a smooth map $(a, b) \mapsto a * b$ defined on an open set
  $D \subseteq E \times E$, with an element $1$ such that $(1, 1) \in D$, $1 * b = b$ and
  $a * 1 = a$ whenever defined, and $(a * b) * c = a * (b * c)$ whenever both sides are
  defined. For $x \in E$, the \emph{left invariant field} with value $x$ at $1$ is
  $\tilde x_l(z) = d(b \mapsto z * b)(1)\,x$ \cite[(5.4)]{GN26}, and the \emph{bracket} is the
  local bracket of these fields at the identity,
  \[
    [x, y] = [\tilde x_l, \tilde y_l](1) = d\tilde y_l(1)\,x - d\tilde x_l(1)\,y .
  \]
\end{definition}

\begin{lemma}[Left invariance]
  \label{lem:xl-mul}
  \lean{LocallyConvex.LieGroups.LocalGroupLaw.xl_mul,
    LocallyConvex.LieGroups.LocalGroupLaw.xl_one}
  \leanok
  \uses{def:local-group-law}
  $\tilde x_l(1) = x$, and $\tilde x_l(a * w) = d(b \mapsto a * b)(w)\,\tilde x_l(w)$
  whenever $a * w$, $w * 1$ and $(a * w) * 1$ are defined.
\end{lemma}

\begin{proof}
  \leanok
  Differentiate $(a * w) * c = a * (w * c)$ in $c$ at $c = 1$ in the direction $x$, with the
  chain rule.
\end{proof}

\begin{proposition}[Brackets of left invariant fields]
  \label{prop:local-bracket-xl}
  \lean{LocallyConvex.LieGroups.LocalGroupLaw.localBracket_xl}
  \leanok
  \uses{lem:xl-mul}
  Near the identity, $[\tilde x_l, \tilde y_l] = \widetilde{[x, y]}_l$ \cite[(5.3)]{GN26}.
\end{proposition}

\begin{proof}
  \leanok
  By Lemma~\ref{lem:xl-mul}, the left translation $b \mapsto a * b$ relates $\tilde x_l$ and
  $\tilde y_l$ to themselves. Local brackets of related fields are related (the second
  derivatives of the translation cancel by Schwarz's theorem; \cite{LCM},
  \texttt{localBracket\_transition}), so
  $[\tilde x_l, \tilde y_l](a) = d(b \mapsto a * b)(1)\,[x, y] = \widetilde{[x, y]}_l(a)$.
\end{proof}

\begin{theorem}[The bracket is a continuous Lie bracket]
  \label{thm:local-lie}
  \lean{LocallyConvex.LieGroups.LocalGroupLaw.bracket_jacobi,
    LocallyConvex.LieGroups.LocalGroupLaw.bracket_swap,
    LocallyConvex.LieGroups.LocalGroupLaw.add_bracket,
    LocallyConvex.LieGroups.LocalGroupLaw.bracket_smul,
    LocallyConvex.LieGroups.LocalGroupLaw.continuous_bracket}
  \leanok
  \uses{prop:local-bracket-xl}
  The bracket of a local group law is bilinear, alternating, continuous, and satisfies the
  Jacobi identity.
\end{theorem}

\begin{proof}
  \leanok
  Bilinearity: $x \mapsto \tilde x_l$ is linear and differentials are linear. Continuity:
  $d\tilde y_l(1)\,x = d^2 m\bigl((1, 1)\bigr)\bigl((0, y), (x, 0)\bigr)$ is jointly continuous
  in $(x, y)$. Jacobi: by Proposition~\ref{prop:local-bracket-xl},
  $[x, [y, z]] = [\tilde x_l, [\tilde y_l, \tilde z_l]](1)$, and local brackets of vector fields
  satisfy the Jacobi identity (\cite{LCM}, \texttt{localBracket\_jacobi}).
\end{proof}

\begin{proposition}[Local homomorphisms]
  \label{prop:local-hom}
  \lean{LocallyConvex.LieGroups.LocalGroupLaw.dMB_bracket_of_hom}
  \leanok
  \uses{prop:local-bracket-xl}
  Let $f$ be $C^2$ near $1$, with $f(1) = 1'$ and $f(a * b) = f(a) *' f(b)$ whenever both
  sides are defined. Then $df(1)[x, y] = [df(1)\,x, df(1)\,y]$.
\end{proposition}

\begin{proof}
  \leanok
  Differentiating $f(z * b) = f(z) *' f(b)$ in $b$ at $1$ shows that $f$ relates $\tilde x_l$
  and $\widetilde{(df(1)x)}_l$; brackets of related fields are related.
\end{proof}

\begin{proposition}[Commutative local group laws]
  \label{prop:local-comm}
  \lean{LocallyConvex.LieGroups.LocalGroupLaw.bracket_eq_zero_of_comm}
  \leanok
  \uses{def:local-group-law}
  If $a * b = b * a$ whenever defined, the bracket vanishes.
\end{proposition}

\begin{proof}
  \leanok
  $d\tilde y_l(1)\,x = d^2 m((1, 1))((0, y), (x, 0))$; commutativity and Schwarz's theorem make
  this symmetric in $x$ and $y$.
\end{proof}

\section{The Lie functor}

\begin{definition}[The Lie algebra of a Lie group]
  \label{def:lie-algebra}
  \lean{LocallyConvex.LieGroups.MBLieGroup.localGroupLaw,
    LocallyConvex.LieGroups.MBLieGroup.LieAlg}
  \leanok
  \uses{def:lie-group, def:local-group-law}
  Let $\varphi$ be the chart of $G$ at the identity. Then
  $a * b = \varphi(\varphi^{-1}(a)\,\varphi^{-1}(b))$ is a local group law on $E$
  \cite[Remark 5.1.8]{GN26}, and $\mathbf{L}(G)$ is $E \cong T_1 G$ with its bracket
  \cite[Definition 5.1.6]{GN26}.
\end{definition}

\begin{theorem}[$\mathbf{L}(G)$ is a locally convex Lie algebra]
  \label{thm:lie-algebra}
  \lean{LocallyConvex.LieGroups.MBLieGroup.LieAlg.instLieRing,
    LocallyConvex.LieGroups.MBLieGroup.LieAlg.instLieAlgebraReal,
    LocallyConvex.LieGroups.MBLieGroup.LieAlg.continuous_lie}
  \leanok
  \uses{def:lie-algebra, thm:local-lie}
  $\mathbf{L}(G)$ is a real Lie algebra with a continuous bracket \cite[Definition 5.1.6]{GN26}.
\end{theorem}

\begin{proof}
  \leanok
  Theorem~\ref{thm:local-lie}.
\end{proof}

\begin{theorem}[The Lie functor]
  \label{thm:lie-functor}
  \lean{LocallyConvex.LieGroups.MBLieGroup.IsHom.lieHom,
    LocallyConvex.LieGroups.MBLieGroup.IsHom.continuous_lieHom,
    LocallyConvex.LieGroups.MBLieGroup.IsHom.lieHom_id,
    LocallyConvex.LieGroups.MBLieGroup.IsHom.lieHom_comp}
  \leanok
  \uses{thm:lie-algebra, def:lie-hom, prop:local-hom}
  For a morphism of Lie groups $\varphi\colon G \to H$, $\mathbf{L}(\varphi) = T_1(\varphi)$ is a
  continuous homomorphism of Lie algebras \cite[Proposition 5.1.10]{GN26}, and
  $\mathbf{L}(\mathrm{id}) = \mathrm{id}$, $\mathbf{L}(\psi \circ \varphi) = \mathbf{L}(\psi)
  \circ \mathbf{L}(\varphi)$ \cite[Corollary 5.1.11]{GN26}.
\end{theorem}

\begin{proof}
  \leanok
  In the charts at the identities, $\varphi$ is a local homomorphism of the local group laws;
  apply Proposition~\ref{prop:local-hom}. Functoriality is the chain rule.
\end{proof}


\begin{theorem}[Left invariant vector fields]
  \label{thm:left-fields}
  \lean{LocallyConvex.LieGroups.MBLieGroup.leftField,
    LocallyConvex.LieGroups.MBLieGroup.tangentMap_leftField,
    LocallyConvex.LieGroups.MBLieGroup.lieBracket_leftField,
    LocallyConvex.LieGroups.MBLieGroup.lieBracket_leftField_one}
  \leanok
  \uses{thm:lie-algebra}
  For $x \in \mathbf{L}(G)$, the vector field $x_l(g) = T_1(\lambda_g)x$ is left invariant,
  $T(\lambda_h) x_l(g) = x_l(hg)$ \cite[(5.2)]{GN26}, and $[x_l, y_l] = [x, y]_l$ for the Lie
  bracket of vector fields \cite[(5.3)]{GN26}. In particular $[x_l, y_l](1) = [x, y]$: the bracket
  of $\mathbf{L}(G)$ is the bracket of left invariant vector fields at the identity
  \cite[Definition 5.1.6]{GN26}.
\end{theorem}

\begin{proof}
  \leanok
  In the chart $\psi$ at $g$, $x_l$ reads $w \mapsto d(b \mapsto \psi(\psi^{-1}(w)\varphi^{-1}(b)))
  (\varphi(1))\,x$. By associativity, the left translation $\tau(b) = \psi(g\varphi^{-1}(b))$
  relates the fields $\tilde x_l$ of the chart $\varphi$ at the identity to these local
  representatives, so local brackets of related fields are related (\cite{LCM},
  \texttt{localBracket\_transition}):
  $[x_l, y_l](g) = d\tau(\varphi 1)[x, y] = [x, y]_l(g)$ \cite[Lemma 3.4.17]{GN26}.
\end{proof}

\begin{proposition}[Abelian Lie groups]
  \label{prop:lie-abelian}
  \lean{LocallyConvex.LieGroups.MBLieGroup.LieAlg.lie_eq_zero_of_comm}
  \leanok
  \uses{thm:lie-algebra, prop:local-comm}
  The Lie algebra of an abelian Lie group is abelian \cite[Example 5.1.9(a)]{GN26}.
\end{proposition}

\begin{proof}
  \leanok
  Proposition~\ref{prop:local-comm}.
\end{proof}

\section{The Lie algebra of \texorpdfstring{$\mathrm{Diff}(K)$}{Diff(K)}}

\begin{lemma}[Evaluating the bracket]
  \label{lem:clm-xl}
  \lean{LocallyConvex.LieGroups.LocalGroupLaw.clm_dMB_xl}
  \leanok
  \uses{def:local-group-law}
  For a continuous linear map $\mathrm{ev}$,
  $\mathrm{ev}(d\tilde y_l(1)\,x) = \partial_t \partial_s\, \mathrm{ev}\bigl((1 + sx) * (1 + ty)\bigr)$
  at $s = t = 0$.
\end{lemma}

\begin{proof}
  \leanok
  $d\tilde y_l(1)\,x$ is a second derivative of the multiplication; $\mathrm{ev}$ commutes with
  derivatives, and Schwarz's theorem exchanges $\partial_s$ and $\partial_t$.
\end{proof}

\begin{theorem}[$\mathbf{L}(\mathrm{Diff}(K)) = \mathcal{V}(K)^{\mathrm{op}}$]
  \label{thm:lie-diff}
  \lean{LocallyConvex.LieGroups.Diffeo.law_mul_apply,
    LocallyConvex.LieGroups.Diffeo.evalAt_dMB_xl,
    LocallyConvex.LieGroups.Diffeo.evalAt_lie,
    LocallyConvex.LieGroups.Diffeo.toField_lie}
  \leanok
  \uses{thm:diff-lie, thm:lie-algebra, lem:clm-xl}
  Let $K$ be a compact manifold with a normalized local addition $\Sigma$. Identify the model
  space $\Gamma_{\mathrm{id}}$ of $\mathrm{Diff}(K)$ with the space of vector fields on $K$. Then
  the bracket of $\mathbf{L}(\mathrm{Diff}(K))$ is the opposite of the Lie bracket of vector
  fields: $[x, y] = [Y, X]$ \cite[Example 5.4.3]{GN26}.
\end{theorem}

\begin{proof}
  \leanok
  In the chart at the identity, $(a * b)(p) = \theta^{-1}(p, \Sigma(a(q)))$ with
  $q = \Sigma(b(p))$. Evaluating at $p$ and using Lemma~\ref{lem:clm-xl},
  $(d\tilde y_l(0)\,x)(p)$ is the mixed derivative at $0$ of
  $G(s, t) = \Theta_p\bigl(\tilde\Sigma(Q(t), s\,\tilde x(Q(t)))\bigr)$, where
  $Q(t) = \varphi_p \Sigma(p, t\,y(p))$, $\Theta_p$ inverts $\Sigma(p, \cdot)$, and $\tilde\Sigma$,
  $\tilde x$ are read in the chart $\varphi_p$ of $K$ at $p$. Since $\Sigma$ is normalized,
  $\tilde\Sigma(z, 0) = z$ and $\partial_v \tilde\Sigma(z, 0) = \mathrm{id}$ near $\varphi_p(p)$,
  and $d\Theta_p(\varphi_p p) = \mathrm{id}$, so the chain rule gives
  $d^2\Theta_p(x(p), y(p)) + d\tilde x \cdot y(p)$. The first term is symmetric in $x$ and
  $y$, hence $[x, y](p) = d\tilde x\cdot \tilde y - d\tilde y \cdot \tilde x = [Y, X](p)$.
  Glöckner and Neeb obtain the formula from the derived action of $\mathrm{Diff}(K)$ on $K$.
\end{proof}
